\documentclass[11pt,a4paper]{article}

% ---- Packages ----
\usepackage[utf8]{inputenc}
\usepackage[T1]{fontenc}
\usepackage{amsmath,amssymb,amsthm}
\usepackage{enumitem}
\usepackage{geometry}
\usepackage{booktabs}
\usepackage{xcolor}
\usepackage[breakable]{tcolorbox}
\usepackage{titlesec}
\usepackage{hyperref}
\usepackage{parskip}
\usepackage{mdframed}

\geometry{margin=2.5cm}

% ---- Colours ----
\definecolor{solcolor}{RGB}{0,100,60}
\definecolor{blockcolor}{RGB}{0,70,130}
\definecolor{hintcolor}{RGB}{140,90,0}
\definecolor{exbg}{RGB}{245,245,235}
\definecolor{warncolor}{RGB}{180,0,0}

% ---- Solution environment ----
\newtcolorbox{solution}{
  colback=solcolor!5,
  colframe=solcolor!70,
  fonttitle=\bfseries,
  title=Solution,
  breakable
}

% ---- Book example environment ----
\newmdenv[
  backgroundcolor=exbg,
  linecolor=blockcolor!40,
  linewidth=1pt,
  roundcorner=4pt,
  innertopmargin=10pt,
  innerbottommargin=10pt,
  innerleftmargin=12pt,
  innerrightmargin=12pt,
  skipabove=8pt,
  skipbelow=8pt
]{bookexample}

% ---- Section formatting ----
\titleformat{\section}{\large\bfseries\color{blockcolor}}{}{0em}{}[\vspace{-0.3em}\rule{\textwidth}{0.4pt}]

% ---- Header ----
\title{\textbf{Tutorial 6 --- Variance of the OLS Estimator\\ and Hypothesis Testing}}
\author{Introductory Econometrics --- TA Session}
\date{Duration: 50 minutes \\ \medskip \small Based on Wooldridge, Theorem 2.2, Sections 2.5 and 4.1--4.2}

\begin{document}
\maketitle
\thispagestyle{empty}

%=============================================================================
\section{Block 1: Motivation --- Why Isn't Unbiasedness Enough? \textnormal{(3 min)}}
%=============================================================================

\begin{bookexample}
\textbf{\color{blockcolor}The Simple Linear Regression Assumptions}

\medskip
We work with the model $y = \beta_0 + \beta_1 x + u$ under the following assumptions. You need to know these by name --- each one will be explicitly cited in the derivations below.

\medskip
\renewcommand{\arraystretch}{1.4}
\begin{tabular}{lll}
\toprule
\textbf{Label} & \textbf{Name} & \textbf{Statement} \\
\midrule
\textbf{SLR.1} & Linear in Parameters & $y = \beta_0 + \beta_1 x + u$ \\
\textbf{SLR.2} & Random Sampling & $\{(x_i, y_i)\}_{i=1}^n$ are i.i.d.\ draws \\
\textbf{SLR.3} & Sample Variation in $x$ & $\text{SST}_x = \sum (x_i - \bar{x})^2 > 0$ \\
\textbf{SLR.4} & Zero Conditional Mean & $E[u \mid x] = 0$ \\
\textbf{SLR.5} & Homoskedasticity & $\text{Var}(u \mid x) = \sigma^2$ (constant) \\
\textbf{SLR.6} & Normality & $u \mid x \sim N(0, \sigma^2)$ \\
\bottomrule
\end{tabular}
\renewcommand{\arraystretch}{1.0}

\medskip
\textbf{What we proved so far:} Under \textbf{SLR.1--SLR.4}, the OLS estimator is \textbf{unbiased}: $E[\hat{\beta}_1] = \beta_1$. But unbiasedness says only that $\hat{\beta}_1$ is centred at $\beta_1$ \emph{on average} across all possible samples. It says nothing about how far any single estimate might be from the truth.

\medskip
\textbf{What we need now:} To know how \emph{precise} $\hat{\beta}_1$ is, we need its variance. Adding \textbf{SLR.5} allows us to derive $\text{Var}(\hat{\beta}_1)$. Adding \textbf{SLR.6} on top gives us the exact sampling distribution, which enables hypothesis tests and confidence intervals.
\end{bookexample}

\begin{bookexample}
\textbf{\color{blockcolor}Theorem 2.2 (Wooldridge) --- Sampling Variance of $\hat{\beta}_1$}

\medskip
Under \textbf{SLR.1--SLR.5}:
\[
  \boxed{\text{Var}(\hat{\beta}_1 \mid \mathbf{x}) = \frac{\sigma^2}{\text{SST}_x}}
  \qquad\text{where}\quad
  \text{SST}_x = \sum_{i=1}^{n}(x_i - \bar{x})^2
\]
The proof relies on a key representation of $\hat{\beta}_1$ that isolates the source of randomness. We derive it step by step.

\medskip
\textbf{Step 1: Start from the OLS formula.}
\[
  \hat{\beta}_1 = \frac{\sum_{i=1}^{n}(x_i - \bar{x})(y_i - \bar{y})}{\sum_{i=1}^{n}(x_i - \bar{x})^2}
  = \frac{\sum(x_i - \bar{x})(y_i - \bar{y})}{\text{SST}_x}
\]

\textbf{Step 2: Substitute the model (\textbf{SLR.1}).} Since $y_i = \beta_0 + \beta_1 x_i + u_i$, taking the sample mean gives $\bar{y} = \beta_0 + \beta_1 \bar{x} + \bar{u}$. Subtracting:
\[
  y_i - \bar{y} = \beta_1(x_i - \bar{x}) + (u_i - \bar{u})
\]

\textbf{Step 3: Expand the numerator.}
\begin{align*}
  \sum(x_i - \bar{x})(y_i - \bar{y})
  &= \sum(x_i - \bar{x})\bigl[\beta_1(x_i - \bar{x}) + (u_i - \bar{u})\bigr] \\
  &= \beta_1 \underbrace{\sum(x_i - \bar{x})^2}_{= \,\text{SST}_x} + \sum(x_i - \bar{x})(u_i - \bar{u})
\end{align*}

\textbf{Step 4: Eliminate $\bar{u}$ from the second term.}
\[
  \sum(x_i - \bar{x})(u_i - \bar{u})
  = \sum(x_i - \bar{x})\,u_i \;-\; \bar{u}\underbrace{\sum(x_i - \bar{x})}_{= \,0}
  = \sum(x_i - \bar{x})\,u_i
\]
The key fact is $\sum(x_i - \bar{x}) = 0$ (deviations from the mean always sum to zero).

\textbf{Step 5: Divide by $\text{SST}_x$.}
\[
  \hat{\beta}_1
  = \frac{\beta_1 \cdot \text{SST}_x + \sum(x_i - \bar{x})\,u_i}{\text{SST}_x}
  = \beta_1 + \frac{\sum(x_i - \bar{x})\,u_i}{\text{SST}_x}
\]
\[
  \boxed{\hat{\beta}_1 = \beta_1 + \sum_{i=1}^{n} w_i u_i,
  \qquad w_i = \frac{x_i - \bar{x}}{\text{SST}_x}}
\]

\textbf{Interpretation:} $\hat{\beta}_1$ equals the true parameter $\beta_1$ plus a weighted sum of the unobserved errors $u_1, \ldots, u_n$. The weights $w_i$ depend only on the $x$-values, so conditional on $\mathbf{x}$ they are constants. The \emph{only} source of randomness in $\hat{\beta}_1$ is the errors $u_i$.
\end{bookexample}

%=============================================================================
\section{Block 2: Deriving and Computing $\text{Var}(\hat{\beta}_1)$ \textnormal{(12 min)}}
%=============================================================================

\textbf{Question 1} (Derivation and numerical computation)

\begin{enumerate}[label=(\alph*)]
\item Starting from the representation $\hat{\beta}_1 = \beta_1 + \sum_{i=1}^{n} w_i u_i$, derive $\text{Var}(\hat{\beta}_1 \mid \mathbf{x}) = \sigma^2/\text{SST}_x$.

\textcolor{hintcolor}{\emph{Hint: Follow these steps --- (i)~Why can you drop $\beta_1$ from the variance? (ii)~Expand $\text{Var}(\sum w_i u_i \mid \mathbf{x})$. Which assumption eliminates the covariance terms? (iii)~Which assumption makes $\text{Var}(u_i \mid \mathbf{x})$ the same for all $i$? (iv)~Simplify $\sum w_i^2$.}}

\begin{solution}
We start from:
\[
  \hat{\beta}_1 = \beta_1 + \sum_{i=1}^{n} w_i u_i
  \qquad\text{where}\quad
  w_i = \frac{x_i - \bar{x}}{\text{SST}_x}
\]

\textbf{Step (i): Drop the constant $\beta_1$.}

By \textbf{SLR.1} (Linear in Parameters), the model is $y = \beta_0 + \beta_1 x + u$, where $\beta_0$ and $\beta_1$ are \emph{fixed, unknown population parameters} --- they are constants, not random variables. Since adding a constant to a random variable does not change its variance (recall: $\text{Var}(a + X) = \text{Var}(X)$ for any constant $a$):
\[
  \text{Var}(\hat{\beta}_1 \mid \mathbf{x})
  = \text{Var}\!\left(\beta_1 + \sum_{i=1}^n w_i u_i \;\middle|\; \mathbf{x}\right)
  = \text{Var}\!\left(\sum_{i=1}^n w_i u_i \;\middle|\; \mathbf{x}\right)
\]
Note also that we condition on $\mathbf{x} = (x_1, \ldots, x_n)$, so the weights $w_i = (x_i - \bar{x})/\text{SST}_x$ are treated as \textbf{constants} (they depend only on the $x$-values). The \textbf{only random variables} in this expression are $u_1, \ldots, u_n$.

\bigskip

\textbf{Step (ii): Expand the variance of the weighted sum.}

By the general formula for the variance of a linear combination:
\[
  \text{Var}\!\left(\sum_{i=1}^n w_i u_i \;\middle|\; \mathbf{x}\right)
  = \underbrace{\sum_{i=1}^n w_i^2\,\text{Var}(u_i \mid \mathbf{x})}_{\text{variance terms}}
  + \underbrace{\sum_{\substack{i,j=1 \\ i \ne j}}^{n} w_i w_j\,\text{Cov}(u_i, u_j \mid \mathbf{x})}_{\text{covariance terms}}
\]

Now we apply \textbf{SLR.2} (Random Sampling): the observations $\{(x_i, y_i)\}_{i=1}^n$ are drawn \emph{independently} from the population. Since $y_i = \beta_0 + \beta_1 x_i + u_i$, the errors $u_1, \ldots, u_n$ are also \textbf{independent of each other} conditional on $\mathbf{x}$.

Independence implies that $\text{Cov}(u_i, u_j \mid \mathbf{x}) = 0$ for all $i \ne j$. Therefore, \textbf{all cross-terms vanish}:
\[
  = \sum_{i=1}^n w_i^2\,\text{Var}(u_i \mid \mathbf{x}) + 0
  = \sum_{i=1}^n w_i^2\,\text{Var}(u_i \mid \mathbf{x})
\]

\bigskip

\textbf{Step (iii): Apply homoskedasticity.}

Now apply \textbf{SLR.5} (Homoskedasticity): the variance of the error term is the \emph{same} for all observations, regardless of the value of $x$:
\[
  \text{Var}(u_i \mid \mathbf{x}) = \sigma^2 \quad\text{for all } i = 1, \ldots, n
\]
Since $\sigma^2$ does not depend on $i$, it can be \textbf{factored out} of the sum:
\[
  \sum_{i=1}^n w_i^2\,\text{Var}(u_i \mid \mathbf{x})
  = \sum_{i=1}^n w_i^2 \cdot \sigma^2
  = \sigma^2 \sum_{i=1}^n w_i^2
\]

\textcolor{warncolor}{\textbf{Why SLR.5 is critical:}} Without homoskedasticity, each observation could have a different error variance $\text{Var}(u_i \mid x_i) = \sigma_i^2$, and we could \emph{not} factor out a single $\sigma^2$. The formula would become $\sum w_i^2 \sigma_i^2$, which depends on each individual $\sigma_i^2$ and is much harder to work with. This is exactly the complication that arises under \textbf{heteroskedasticity}.

\bigskip

\textbf{Step (iv): Simplify $\sum w_i^2$.}

Recall that $w_i = (x_i - \bar{x})/\text{SST}_x$. Squaring and summing:
\[
  \sum_{i=1}^n w_i^2
  = \sum_{i=1}^n \frac{(x_i - \bar{x})^2}{\text{SST}_x^2}
  = \frac{1}{\text{SST}_x^2} \sum_{i=1}^n (x_i - \bar{x})^2
  = \frac{\text{SST}_x}{\text{SST}_x^2}
  = \frac{1}{\text{SST}_x}
\]
In the third equality, we used the definition $\text{SST}_x = \sum_{i=1}^n (x_i - \bar{x})^2$. Note that \textbf{SLR.3} (Sample Variation in $x$) guarantees $\text{SST}_x > 0$, so this division is valid.

\bigskip

\textbf{Combining all four steps:}
\[
  \text{Var}(\hat{\beta}_1 \mid \mathbf{x})
  \underset{\text{(i)}}{=} \text{Var}\!\left(\sum w_i u_i \mid \mathbf{x}\right)
  \underset{\text{(ii)}}{=} \sum w_i^2\,\text{Var}(u_i \mid \mathbf{x})
  \underset{\text{(iii)}}{=} \sigma^2 \sum w_i^2
  \underset{\text{(iv)}}{=} \sigma^2 \cdot \frac{1}{\text{SST}_x}
\]
\[
  \boxed{\text{Var}(\hat{\beta}_1 \mid \mathbf{x}) = \frac{\sigma^2}{\text{SST}_x}} \qquad\square
\]

\medskip
\textbf{Summary of assumptions used:}
\begin{itemize}[itemsep=2pt]
  \item \textbf{SLR.1} (Linear in Parameters): $\beta_1$ is a constant, so it drops from the variance.
  \item \textbf{SLR.2} (Random Sampling): errors are independent, so covariance terms are zero.
  \item \textbf{SLR.3} (Sample Variation): $\text{SST}_x > 0$, so division is valid.
  \item \textbf{SLR.5} (Homoskedasticity): all $\text{Var}(u_i \mid \mathbf{x}) = \sigma^2$, so $\sigma^2$ factors out.
\end{itemize}
Note: SLR.4 (Zero Conditional Mean) is \emph{not} needed for the variance formula --- it was needed for unbiasedness but not here.
\end{solution}

\item Suppose $n = 5$ observations with $x$-values $\{1, 3, 5, 7, 9\}$ and $\sigma^2 = 10$. Compute $\bar{x}$, $\text{SST}_x$, $\text{Var}(\hat{\beta}_1)$, and $\text{sd}(\hat{\beta}_1)$.

\begin{solution}
\textbf{Step 1: Compute $\bar{x}$.}
\[
  \bar{x} = \frac{1}{n}\sum_{i=1}^n x_i = \frac{1 + 3 + 5 + 7 + 9}{5} = \frac{25}{5} = 5
\]

\textbf{Step 2: Compute each deviation $(x_i - \bar{x})$ and its square.}

\medskip
\begin{center}
\begin{tabular}{ccc}
\toprule
$x_i$ & $x_i - \bar{x}$ & $(x_i - \bar{x})^2$ \\
\midrule
1 & $1 - 5 = -4$ & $(-4)^2 = 16$ \\
3 & $3 - 5 = -2$ & $(-2)^2 = 4$ \\
5 & $5 - 5 = 0$  & $0^2 = 0$ \\
7 & $7 - 5 = 2$  & $2^2 = 4$ \\
9 & $9 - 5 = 4$  & $4^2 = 16$ \\
\midrule
  &  & $\text{SST}_x = 40$ \\
\bottomrule
\end{tabular}
\end{center}

\textbf{Step 3: Apply the formula.}
\[
  \text{Var}(\hat{\beta}_1)
  = \frac{\sigma^2}{\text{SST}_x}
  = \frac{10}{40}
  = \boxed{0.25}
\]
\[
  \text{sd}(\hat{\beta}_1)
  = \sqrt{\text{Var}(\hat{\beta}_1)}
  = \sqrt{0.25}
  = \boxed{0.5}
\]

\textbf{Interpretation:} Across repeated samples with these same $x$-values, the OLS slope estimate $\hat{\beta}_1$ would have a standard deviation of $0.5$ --- on average, our estimates will be about $0.5$ units away from the true $\beta_1$.
\end{solution}

\item Based on the formula $\text{Var}(\hat{\beta}_1) = \sigma^2/\text{SST}_x$, explain what happens to the precision of $\hat{\beta}_1$ when:
\begin{enumerate}[label=(\roman*)]
  \item the error variance $\sigma^2$ increases;
  \item we add more observations that are spread out (not concentrated at $\bar{x}$);
  \item we add observations that are all equal to $\bar{x}$.
\end{enumerate}

\begin{solution}
The formula $\text{Var}(\hat{\beta}_1) = \sigma^2/\text{SST}_x$ has two ingredients: $\sigma^2$ in the numerator and $\text{SST}_x$ in the denominator.

\begin{enumerate}[label=(\roman*)]
  \item \textbf{$\sigma^2$ increases $\Longrightarrow$ Var increases (less precise).}

  $\sigma^2$ measures the noise in the error term. More noise means the data points are more scattered around the regression line, making it harder to pin down the slope.

  \emph{Numerical example:} With our data ($\text{SST}_x = 40$), doubling $\sigma^2$ from 10 to 20 doubles $\text{Var}(\hat{\beta}_1)$ from $0.25$ to $0.50$.

  \item \textbf{Adding spread-out observations increases $\text{SST}_x$ $\Longrightarrow$ Var decreases (more precise).}

  Each new observation at $x_i \ne \bar{x}$ contributes $(x_i - \bar{x})^2 > 0$ to $\text{SST}_x$. A larger $\text{SST}_x$ in the denominator shrinks $\text{Var}(\hat{\beta}_1)$.

  \emph{Numerical example:} Adding $x_6 = 11$ (far from the mean) would add $(11 - \bar{x}_{\text{new}})^2$ to $\text{SST}_x$, substantially reducing the variance.

  \textbf{Intuition:} More variation in $x$ gives us more ``leverage'' to estimate the slope. Imagine fitting a line through data points that are all bunched together vs.\ spread far apart --- the spread-out data pins the slope down much more precisely.

  \item \textbf{Adding observations at $x = \bar{x}$ does \emph{not} help.}

  If $x_i = \bar{x}$, then $(x_i - \bar{x})^2 = 0$, so these observations contribute \textbf{nothing} to $\text{SST}_x$. $\text{Var}(\hat{\beta}_1)$ is unchanged.

  \emph{Numerical example:} Adding 100 observations all at $x = 5$ (the mean) keeps $\text{SST}_x = 40$, so $\text{Var}(\hat{\beta}_1)$ stays at $0.25$.

  \textbf{Intuition:} To estimate a slope, you need to see how $y$ changes as $x$ changes. If all your new data have the same $x$, you learn nothing new about the slope --- you only learn more about the intercept.
\end{enumerate}
\end{solution}

\end{enumerate}

%=============================================================================
\section{Block 3: From $\sigma^2$ to Standard Errors \textnormal{(10 min)}}
%=============================================================================

\begin{bookexample}
\textbf{\color{blockcolor}Estimating $\sigma^2$ and the $t$-Distribution}

\smallskip
The formula $\text{Var}(\hat{\beta}_1) = \sigma^2/\text{SST}_x$ involves $\sigma^2 = \text{Var}(u \mid \mathbf{x})$, which is \textbf{unknown} (we never observe the true errors $u_i$). We estimate it from the \textbf{residuals} $\hat{u}_i = y_i - \hat{y}_i$:
\[
  \hat{\sigma}^2 = \frac{1}{n-2}\sum_{i=1}^{n} \hat{u}_i^2
  = \frac{\text{SSR}}{n-2}
\]
\textbf{Why $n-2$ and not $n$?} We estimated two parameters ($\hat{\beta}_0$ and $\hat{\beta}_1$) to compute the residuals. This ``uses up'' 2 degrees of freedom. Dividing by $n-2$ corrects for this and makes $\hat{\sigma}^2$ an \textbf{unbiased} estimator of $\sigma^2$: $E[\hat{\sigma}^2] = \sigma^2$.

\smallskip
The \textbf{standard error} of $\hat{\beta}_1$ replaces $\sigma$ with $\hat{\sigma}$:
\[
  \text{SE}(\hat{\beta}_1) = \frac{\hat{\sigma}}{\sqrt{\text{SST}_x}}
\]
In R, $\hat{\sigma}$ is reported as ``\texttt{Residual standard error}'' and $\text{SE}(\hat{\beta}_1)$ appears in the ``\texttt{Std.\ Error}'' column.

\smallskip
Adding \textbf{SLR.6} (Normality: $u \mid x \sim N(0, \sigma^2)$), the $t$-statistic
\[
  t = \frac{\hat{\beta}_1 - \beta_1}{\text{SE}(\hat{\beta}_1)} \sim t_{n-2}
\]
follows a \textbf{$t$-distribution} with $n - 2$ degrees of freedom. The $t$ (not $N(0,1)$) arises because $\hat{\sigma}$ in the denominator is \emph{itself a random variable} --- it varies from sample to sample, adding extra uncertainty beyond the randomness of $\hat{\beta}_1$.

\smallskip
\textbf{Key distinction:}
\begin{itemize}[leftmargin=1.5em, itemsep=2pt]
  \item If $\sigma$ were \textbf{known}: $Z = \dfrac{\hat{\beta}_1 - \beta_1}{\sigma/\sqrt{\text{SST}_x}} \sim N(0,1)$ \quad (standard normal)
  \item Since $\sigma$ is \textbf{unknown}: $t = \dfrac{\hat{\beta}_1 - \beta_1}{\hat{\sigma}/\sqrt{\text{SST}_x}} \sim t_{n-2}$ \quad ($t$-distribution, heavier tails)
\end{itemize}
As $n \to \infty$, $\hat{\sigma} \to \sigma$ and $t_{n-2} \to N(0,1)$, so the distinction vanishes in large samples.
\end{bookexample}

\bigskip

The following R output will be used for Questions~2 and~3. An econometrician regressed \verb!y! on \verb!x! using 22 observations.
\small
\begin{verbatim}
> model <- lm(y ~ x, data = mydata)
> summary(model)
...
Coefficients:
            Estimate Std. Error t value Pr(>|t|)
(Intercept)   5.2000     1.3000   4.000 0.000742 ***
x             2.4000     0.8000       A        B
---
Signif. codes:  0 '***' 0.001 '**' 0.01 '*' 0.05 '.' 0.1 ' ' 1

Residual standard error: 4 on 20 degrees of freedom

> confint(model)
                2.5 %   97.5 %
(Intercept)  2.488248 7.911752
x                   C        D
\end{verbatim}
\normalsize

\bigskip
\textbf{Question 2} (Reading R output --- variance and standard errors)

\begin{enumerate}[label=(\alph*)]
\item From the R output, identify $\hat{\sigma}$ (the residual standard error) and the number of observations $n$. Then compute $\hat{\sigma}^2$ and $\text{SST}_x$.

\textcolor{hintcolor}{\emph{Hint: Use $\text{SE}(\hat{\beta}_1) = \hat{\sigma}/\sqrt{\text{SST}_x}$ and solve for $\text{SST}_x$.}}

\begin{solution}
\textbf{Reading $\hat{\sigma}$ from the output:} The line ``\texttt{Residual standard error: 4 on 20 degrees of freedom}'' tells us two things:
\begin{itemize}
  \item $\hat{\sigma} = 4$ (the estimated standard deviation of the errors)
  \item Degrees of freedom $= n - 2 = 20$, which means $n = 22$ observations.
\end{itemize}

\textbf{Computing $\hat{\sigma}^2$:}
\[
  \hat{\sigma}^2 = \hat{\sigma}^{\,2} = 4^2 = \boxed{16}
\]

\textbf{Computing $\text{SST}_x$:} From the output, the ``\texttt{Std.\ Error}'' column for \verb!x! gives $\text{SE}(\hat{\beta}_1) = 0.8$. Using the formula $\text{SE}(\hat{\beta}_1) = \hat{\sigma}/\sqrt{\text{SST}_x}$:
\[
  0.8 = \frac{4}{\sqrt{\text{SST}_x}}
\]
Solve for $\sqrt{\text{SST}_x}$: multiply both sides by $\sqrt{\text{SST}_x}$ and divide by $0.8$:
\[
  \sqrt{\text{SST}_x} = \frac{4}{0.8} = 5
\]
Square both sides:
\[
  \text{SST}_x = 5^2 = \boxed{25}
\]

\textbf{Verification:} $\text{SE}(\hat{\beta}_1) = \hat{\sigma}/\sqrt{\text{SST}_x} = 4/\sqrt{25} = 4/5 = 0.8$ \checkmark
\end{solution}

\item Explain in one or two sentences why replacing $\sigma$ with $\hat{\sigma}$ changes the distribution of the test statistic from standard normal to $t$.

\begin{solution}
When $\sigma$ is \textbf{known}, $(\hat{\beta}_1 - \beta_1)/(\sigma/\sqrt{\text{SST}_x})$ is a ratio of a normal random variable (the numerator) over a \emph{constant} (the denominator), so the ratio is exactly $N(0,1)$.

When we \textbf{replace} $\sigma$ with $\hat{\sigma}$, the denominator $\hat{\sigma}/\sqrt{\text{SST}_x}$ becomes a \emph{random variable} --- it fluctuates from sample to sample because $\hat{\sigma}$ is computed from the data. The ratio is now a normal random variable divided by a random estimate of its standard deviation. This ratio has \textbf{heavier tails} than the normal (sometimes $\hat{\sigma}$ underestimates $\sigma$, inflating $|t|$), producing the $t_{n-2}$ distribution.

\textbf{Intuition:} The $t$-distribution is ``wider'' than the normal to account for the fact that we \emph{don't know} $\sigma$ exactly. As the sample size grows, $\hat{\sigma}$ becomes a better and better estimate of $\sigma$, the extra uncertainty shrinks, and $t_{n-2} \to N(0,1)$.
\end{solution}

\item Find $A$ (the $t$-value) and $B$ (the $p$-value) in the output. What null hypothesis does this $t$-statistic test?

\begin{solution}
\textbf{What R reports by default:} The $t$-value and $p$-value in R's \verb!summary()! output always test the null hypothesis
\[
  H_0\colon \beta_1 = 0 \quad\text{vs.}\quad H_1\colon \beta_1 \ne 0
\]
That is, R tests whether the coefficient is significantly different from \textbf{zero} (two-sided).

\textbf{Finding $A$ (t-value):} The $t$-statistic for testing $H_0\colon \beta_1 = 0$ is:
\[
  A = t = \frac{\hat{\beta}_1 - 0}{\text{SE}(\hat{\beta}_1)}
  = \frac{\text{Estimate}}{\text{Std.\ Error}}
  = \frac{2.4000}{0.8000}
  = \boxed{3.000}
\]

\textbf{Finding $B$ (p-value):} The $p$-value is the probability of observing a $t$-statistic as extreme as $|t| = 3$ or more, \emph{if $H_0$ were true}:
\[
  B = P(|t_{20}| > 3) = 2 \cdot P(t_{20} > 3)
\]
The factor of 2 appears because this is a \textbf{two-sided} test (we count both tails). In R: \verb!2 * pt(-3, df = 20)! $\approx \boxed{0.007}$.

\textbf{Meaning:} If $\beta_1$ were truly zero, there would be only a $0.7\%$ chance of observing a $t$-statistic as large as $3$ in absolute value. This is strong evidence against $H_0$.
\end{solution}

\end{enumerate}

%=============================================================================
\section{Block 4: Hypothesis Testing \textnormal{(13 min)}}
%=============================================================================

\begin{bookexample}
\textbf{\color{blockcolor}Hypothesis Testing and Confidence Intervals --- The Procedure}

\smallskip
\textbf{Two-sided test} of $H_0\colon \beta_1 = \beta_{1,0}$ vs.\ $H_1\colon \beta_1 \ne \beta_{1,0}$:
\begin{enumerate}[leftmargin=2em,itemsep=2pt]
  \item State the hypotheses ($H_0$ and $H_1$) and significance level $\alpha$.
  \item Compute the test statistic: $t = (\hat{\beta}_1 - \beta_{1,0})/\text{SE}(\hat{\beta}_1)$.
  \item Find the critical value: $t_{\alpha/2,\, n-2}$.
  \item Decision rule: Reject $H_0$ if $|t| > t_{\alpha/2,\, n-2}$.
  \item State the conclusion in context.
\end{enumerate}

\smallskip
\textbf{One-sided test} of $H_0\colon \beta_1 \le \beta_{1,0}$ vs.\ $H_1\colon \beta_1 > \beta_{1,0}$:
\begin{itemize}[leftmargin=1.5em,itemsep=2pt]
  \item Same $t$-statistic.
  \item Reject $H_0$ if $t > t_{\alpha,\, n-2}$ (entire $\alpha$ in one tail).
  \item Since $t_{\alpha,\, n-2} < t_{\alpha/2,\, n-2}$, it is \textbf{easier to reject} in the specified direction.
\end{itemize}

\smallskip
\textbf{$(1-\alpha)$ Confidence interval:}
$\hat{\beta}_1 \pm t_{\alpha/2,\, n-2} \cdot \text{SE}(\hat{\beta}_1)$.

\smallskip
\textbf{CI--test equivalence (two-sided only):} Reject $H_0\colon \beta_1 = \beta_{1,0}$ at level $\alpha$ if and only if $\beta_{1,0}$ falls \textbf{outside} the $(1-\alpha)$ CI.
\end{bookexample}

\medskip
\textit{Use the R output from Block~3. You may use the following critical values for the $t_{20}$ distribution: $t_{0.025,\, 20} = 2.086$ (in R: \texttt{qt(0.975, 20)}) and $t_{0.05,\, 20} = 1.725$ (in R: \texttt{qt(0.95, 20)}).}

\bigskip
\textbf{Question 3} (Hypothesis testing and confidence intervals)

\begin{enumerate}[label=(\alph*)]
\item Using $A$ and $B$ from Question~2, test $H_0\colon \beta_1 = 0$ against $H_1\colon \beta_1 \ne 0$ at the 5\% significance level. State your conclusion.

\begin{solution}
\textbf{Step 1: State hypotheses and significance level.}
\begin{itemize}
  \item $H_0\colon \beta_1 = 0$ (the variable $x$ has no effect on $y$)
  \item $H_1\colon \beta_1 \ne 0$ (the variable $x$ has some effect on $y$)
  \item This is a \textbf{two-sided} test at $\alpha = 0.05$.
\end{itemize}

\textbf{Step 2: Test statistic.}

From Q2(c): $t = A = 3.000$.

\textbf{Step 3: Critical value.}

For a two-sided test at $\alpha = 0.05$ with $\text{df} = 20$: $t_{0.025,\, 20} = 2.086$.

\textbf{Step 4: Decision.}

Reject $H_0$ if $|t| > 2.086$. Since $|t| = 3.000 > 2.086$, we \textbf{reject $H_0$}.

\textbf{Alternative (using the $p$-value):} Reject $H_0$ if $p < \alpha$. Since $p = B = 0.007 < 0.05 = \alpha$, we \textbf{reject $H_0$}. \checkmark

\textbf{Step 5: Conclusion.}

At the 5\% significance level, we reject the null hypothesis that $\beta_1 = 0$. There is statistically significant evidence that $x$ has an effect on $y$.
\end{solution}

\item Find $C$ and $D$ (the 95\% confidence interval for $\beta_1$). Show your work.

\begin{solution}
\textbf{Step 1: Formula.}

The $(1-\alpha) = 95\%$ confidence interval for $\beta_1$ is:
\[
  \hat{\beta}_1 \pm t_{\alpha/2,\, n-2} \cdot \text{SE}(\hat{\beta}_1)
\]

\textbf{Step 2: Plug in values.}
\begin{itemize}
  \item $\hat{\beta}_1 = 2.4$ (from the ``Estimate'' column)
  \item $t_{0.025,\, 20} = 2.086$ (critical value for 95\% CI with 20 df)
  \item $\text{SE}(\hat{\beta}_1) = 0.8$ (from the ``Std.\ Error'' column)
\end{itemize}

\textbf{Step 3: Compute the margin of error.}
\[
  \text{Margin of error} = t_{0.025,\, 20} \times \text{SE}(\hat{\beta}_1)
  = 2.086 \times 0.8 = 1.669
\]

\textbf{Step 4: Compute the bounds.}
\begin{align*}
  C = \text{Lower bound} &= \hat{\beta}_1 - \text{Margin of error} = 2.4 - 1.669 = \boxed{0.731} \\
  D = \text{Upper bound} &= \hat{\beta}_1 + \text{Margin of error} = 2.4 + 1.669 = \boxed{4.069}
\end{align*}

\textbf{Step 5: State the interval.}

The 95\% CI for $\beta_1$ is $[0.731,\; 4.069]$.

\textbf{Interpretation:} We are 95\% confident that the true $\beta_1$ lies between $0.731$ and $4.069$. This means that if we repeated the sampling process many times and computed a 95\% CI each time, approximately 95\% of those intervals would contain the true $\beta_1$.

\textbf{Consistency check with part~(a):} The CI does not contain $0$, which is consistent with our rejection of $H_0\colon \beta_1 = 0$ at 5\%.
\end{solution}

\item Using the confidence interval from~(b), test $H_0\colon \beta_1 = 1$ against $H_1\colon \beta_1 \ne 1$ at the 5\% significance level.

\begin{solution}
\textbf{Step 1: State hypotheses.}
\begin{itemize}
  \item $H_0\colon \beta_1 = 1$
  \item $H_1\colon \beta_1 \ne 1$
  \item Two-sided test at $\alpha = 0.05$.
\end{itemize}

\textbf{Step 2: Apply the CI--test equivalence.}

For a two-sided test, we reject $H_0\colon \beta_1 = \beta_{1,0}$ at level $\alpha$ if and only if the hypothesized value $\beta_{1,0}$ falls \textbf{outside} the $(1-\alpha)$ CI. From part~(b), the 95\% CI is $[0.731,\; 4.069]$.

\textbf{Step 3: Check whether $\beta_{1,0} = 1$ is inside or outside the CI.}

Since $0.731 < 1 < 4.069$, the value $\beta_{1,0} = 1$ is \textbf{inside} the 95\% CI.

\textbf{Step 4: Decision.}

We \textbf{fail to reject} $H_0$ at the 5\% significance level. There is not enough evidence to conclude that $\beta_1 \ne 1$.

\medskip
\textbf{Verification via $t$-test:}
\[
  t = \frac{\hat{\beta}_1 - 1}{\text{SE}(\hat{\beta}_1)} = \frac{2.4 - 1}{0.8} = \frac{1.4}{0.8} = 1.75
\]
Decision: $|t| = 1.75 < 2.086 = t_{0.025,20}$, so we fail to reject. \checkmark

Both methods (CI approach and $t$-test) give the same answer --- they are mathematically equivalent for two-sided tests.
\end{solution}

\item Now test $H_0\colon \beta_1 \le 1$ against $H_1\colon \beta_1 > 1$ at the 5\% significance level.

\begin{solution}
\textbf{Step 1: State hypotheses and significance level.}
\begin{itemize}
  \item $H_0\colon \beta_1 \le 1$
  \item $H_1\colon \beta_1 > 1$
  \item This is a \textbf{one-sided} (right-tailed) test at $\alpha = 0.05$.
\end{itemize}

\textbf{Step 2: Compute the test statistic.}

The $t$-statistic uses the boundary value of $H_0$ (i.e., $\beta_{1,0} = 1$):
\[
  t = \frac{\hat{\beta}_1 - \beta_{1,0}}{\text{SE}(\hat{\beta}_1)}
  = \frac{2.4 - 1}{0.8}
  = \frac{1.4}{0.8}
  = 1.75
\]
This is the same $t$-statistic as in part~(c). The difference will be in the critical value.

\textbf{Step 3: Find the critical value.}

For a \textbf{one-sided} test at $\alpha = 0.05$ with $\text{df} = 20$: $t_{0.05,\, 20} = 1.725$.

Notice this is \emph{smaller} than the two-sided critical value ($2.086$). The one-sided test places the entire 5\% rejection probability in one tail instead of splitting it as 2.5\% in each tail.

\textbf{Step 4: Decision rule and decision.}

For a right-tailed test: Reject $H_0$ if $t > t_{0.05, 20} = 1.725$.

Since $t = 1.75 > 1.725$, we \textbf{reject $H_0$}.

\textbf{Step 5: Conclusion.}

At the 5\% significance level, we reject $H_0\colon \beta_1 \le 1$ in favor of $H_1\colon \beta_1 > 1$. There is statistically significant evidence that $\beta_1 > 1$.
\end{solution}

\item Parts~(c) and~(d) test closely related hypotheses about $\beta_1 = 1$, yet give opposite conclusions. Explain why.

\begin{solution}
\textbf{The key: different critical values.}

\begin{center}
\begin{tabular}{lcccc}
\toprule
\textbf{Test} & $t$-stat & \textbf{Critical value} & \textbf{Exceeds CV?} & \textbf{Decision} \\
\midrule
Two-sided (c) & $1.75$ & $t_{0.025,20} = 2.086$ & $1.75 < 2.086$ & Fail to reject \\
One-sided (d)  & $1.75$ & $t_{0.05,20} = 1.725$  & $1.75 > 1.725$ & Reject \\
\bottomrule
\end{tabular}
\end{center}

\textbf{Why the critical values differ:}
\begin{itemize}
  \item The \textbf{two-sided test} splits the 5\% significance level across \emph{both} tails: 2.5\% in the left tail, 2.5\% in the right tail. The critical value $t_{0.025,20} = 2.086$ is large because it must capture only 2.5\% in each tail.

  \item The \textbf{one-sided test} places the \emph{entire} 5\% in the right tail (the direction of $H_1\colon \beta_1 > 1$). The critical value $t_{0.05,20} = 1.725$ is smaller because it captures 5\% in one tail.
\end{itemize}

\textbf{Since $t = 1.75$ falls between the two critical values} ($1.725 < 1.75 < 2.086$), it is large enough to reject in the one-sided test but not large enough to reject in the two-sided test.

\medskip
\textbf{The trade-off:}
\begin{itemize}
  \item The \textbf{one-sided test} has more \textbf{power} (higher probability of rejecting $H_0$) when the true $\beta_1$ lies in the direction of $H_1$ (here, $\beta_1 > 1$). It achieves this by concentrating all its rejection region on one side.

  \item However, the one-sided test has \textbf{zero power} against departures in the \emph{opposite} direction ($\beta_1 < 1$). If $\beta_1$ were actually much less than 1, the one-sided test would never reject.

  \item The \textbf{two-sided test} can detect departures in \emph{either} direction, but at the cost of lower power in any single direction.
\end{itemize}

\textbf{When to use one-sided:} Only when you have strong theoretical reasons to expect the effect in a specific direction \emph{before seeing the data}. Otherwise, use a two-sided test.
\end{solution}

\end{enumerate}

%=============================================================================
\section{Block 5: Type I/II Errors, Size, and Power \textnormal{(12 min)}}
%=============================================================================

\begin{bookexample}
\textbf{\color{blockcolor}Type I and Type II Errors}

\smallskip
When we perform a hypothesis test, we make a decision (reject or fail to reject $H_0$) based on sample data. Since the data are random, we can make \textbf{mistakes}:

\medskip
\begin{center}
\renewcommand{\arraystretch}{1.4}
\begin{tabular}{l|cc}
 & \textbf{$H_0$ is actually true} & \textbf{$H_0$ is actually false} \\
\hline
\textbf{Reject $H_0$}      & \textcolor{red}{Type I Error}  & Correct (Power) \\
\textbf{Fail to reject $H_0$} & Correct              & \textcolor{red}{Type II Error} \\
\end{tabular}
\renewcommand{\arraystretch}{1.0}
\end{center}

\smallskip
\begin{itemize}[leftmargin=1.5em, itemsep=4pt]
  \item \textbf{Type I Error (false positive):} Rejecting $H_0$ when $H_0$ is \emph{true}. We conclude there is an effect, but in reality there is none.

  \item \textbf{Type II Error (false negative):} Failing to reject $H_0$ when $H_0$ is \emph{false}. We miss a real effect.
\end{itemize}

\medskip
\textbf{Example 1 --- Courtroom trial:}
\begin{itemize}[leftmargin=1.5em, itemsep=2pt]
  \item $H_0$: The defendant is innocent.
  \item $H_1$: The defendant is guilty.
  \item \textbf{Type I:} Convicting an innocent person. The jury finds them guilty, but they didn't commit the crime.
  \item \textbf{Type II:} Acquitting a guilty person. The jury finds them not guilty, but they actually did it.
  \item The ``beyond reasonable doubt'' standard sets a very low $\alpha$ --- society considers convicting an innocent person (Type~I) worse than letting a guilty person go free (Type~II).
\end{itemize}

\medskip
\textbf{Example 2 --- Drug trial:}
\begin{itemize}[leftmargin=1.5em, itemsep=2pt]
  \item $H_0\colon \beta_1 = 0$ (the new drug has no effect on blood pressure).
  \item $H_1\colon \beta_1 \ne 0$ (the drug has an effect).
  \item \textbf{Type I:} The trial concludes the drug works (rejects $H_0$), but in reality it has no effect. A useless drug is marketed.
  \item \textbf{Type II:} The trial fails to find an effect (fails to reject $H_0$), but the drug actually works. A beneficial drug is never brought to market.
  \item $\alpha = 0.05$ means: \emph{if the drug truly has no effect}, there is at most a 5\% chance we mistakenly conclude it works.
\end{itemize}
\end{bookexample}

\begin{bookexample}
\textbf{\color{blockcolor}Size and Power of a Test}

\smallskip
\textbf{Size} $= P(\text{reject } H_0 \mid H_0 \text{ is true}) = P(\text{Type I Error}) = \alpha$.

The size is the probability of making a Type~I error. By construction, a test at significance level $\alpha$ has size $\alpha$: we \emph{choose} the critical value to make this probability exactly $\alpha$.

\medskip
\textbf{Power} $= P(\text{reject } H_0 \mid H_0 \text{ is false}) = 1 - P(\text{Type II Error})$.

The power is the probability of \emph{correctly} detecting a real effect. It depends on the \textbf{true value} $\beta^*$: the farther $\beta^*$ is from $\beta_0$ (the null value), the higher the power.

\medskip
\textbf{What we want:} Low Type~I error ($\alpha$ small) AND high power ($1 - \text{Type II}$ large). But these are in tension: making $\alpha$ smaller (harder to reject) also reduces power. The standard compromise is $\alpha = 0.05$ and hoping for power $\ge 0.80$.
\end{bookexample}

\bigskip

\textbf{Question 4} (Type I/II errors and power --- numerical exercise)

\medskip
\textit{To simplify the calculations, assume $\sigma_{\hat{\beta}}$ is \textbf{known}, so the test statistic is $Z = (\hat{\beta} - \beta_0)/\sigma_{\hat{\beta}} \sim N(0,1)$ under $H_0$ (standard normal, not $t$).}

\medskip
Consider testing $H_0\colon \beta = 0$ vs.\ $H_1\colon \beta \ne 0$ at $\alpha = 0.05$, with $\sigma_{\hat{\beta}} = 2$.

\begin{enumerate}[label=(\alph*)]
\item Define Type~I and Type~II error in the context of this test. Give a concrete interpretation if $\beta$ measures the effect of a job training program on wages (in thousands of dollars).

\begin{solution}
\textbf{Type I Error} $= P(\text{reject } H_0 \mid H_0 \text{ is true})$:

We reject $H_0\colon \beta = 0$ even though $\beta$ is actually zero. In context: we conclude that the job training program raises wages, but in reality the program has \textbf{no effect}. We waste resources implementing a useless program.

\medskip

\textbf{Type II Error} $= P(\text{fail to reject } H_0 \mid H_0 \text{ is false})$:

We fail to reject $H_0\colon \beta = 0$ even though $\beta \ne 0$ (the program actually works). In context: we conclude that the job training program has no effect, and we \textbf{cancel a program that actually helps workers}. The real benefit goes undetected.
\end{solution}

\item The two-sided test rejects $H_0$ when $|Z| > 1.96$. Show that the \textbf{size} of this test is exactly $0.05$.

\begin{solution}
The size is the probability of rejecting $H_0$ \emph{when $H_0$ is true}. Under $H_0\colon \beta = 0$:
\[
  Z = \frac{\hat{\beta} - 0}{\sigma_{\hat{\beta}}} = \frac{\hat{\beta}}{2} \sim N(0, 1)
\]
The test rejects when $|Z| > 1.96$. So:
\begin{align*}
  \text{Size}
  &= P(|Z| > 1.96 \mid H_0)  \\
  &= P(Z > 1.96) + P(Z < -1.96) & &\text{(split into two tails)}\\
  &= (1 - \Phi(1.96)) + \Phi(-1.96) & &\text{(using the normal CDF)} \\
  &= 0.025 + 0.025  & &\text{(by symmetry of the normal)} \\
  &= \boxed{0.05}
\end{align*}
This confirms that the critical value $z_{0.025} = 1.96$ was \emph{chosen} to make the Type~I error probability exactly $\alpha = 0.05$. That is the \emph{definition} of a test at significance level~$\alpha$.
\end{solution}

\item Compute the \textbf{power} of this test when the true value is $\beta^* = 6$ (the program truly raises wages by \$6{,}000).

\textcolor{hintcolor}{\emph{Hint: Under $\beta^* = 6$, $\hat{\beta} \sim N(6, 4)$, so $Z = \hat{\beta}/2 \sim N(3, 1)$. Compute $P(|Z| > 1.96)$ using the substitution $W = Z - 3 \sim N(0,1)$. You may use: $\Phi(1.04) = 0.851$.}}

\begin{solution}
\textbf{Step 1: Distribution of $Z$ under the true $\beta^*$.}

If $\beta^* = 6$, then $\hat{\beta} \sim N(\beta^*, \sigma_{\hat{\beta}}^2) = N(6, 4)$. The test statistic (computed as if $H_0$ were true) is:
\[
  Z = \frac{\hat{\beta} - 0}{2} = \frac{\hat{\beta}}{2} \sim N\!\left(\frac{6}{2},\, 1\right) = N(3, 1)
\]
So under $\beta^* = 6$, $Z$ is not $N(0,1)$ but rather $N(3, 1)$ --- it is shifted to the right by $\delta = \beta^*/\sigma_{\hat{\beta}} = 6/2 = 3$.

\textbf{Step 2: Compute power $= P(|Z| > 1.96)$ under $Z \sim N(3, 1)$.}

Split into two tails:
\[
  \text{Power} = P(Z > 1.96) + P(Z < -1.96)
\]
Substitute $W = Z - 3 \sim N(0,1)$, so $Z = W + 3$:

\textbf{Right tail:}
\[
  P(Z > 1.96) = P(W + 3 > 1.96) = P(W > 1.96 - 3) = P(W > -1.04) = \Phi(1.04) = 0.851
\]

\textbf{Left tail:}
\[
  P(Z < -1.96) = P(W + 3 < -1.96) = P(W < -1.96 - 3) = P(W < -4.96) = \Phi(-4.96) \approx 0.000
\]

\textbf{Combining:}
\[
  \text{Power} = 0.851 + 0.000 = \boxed{0.851}
\]

\textbf{Interpretation:} If the job training program truly raises wages by \$6{,}000 (with $\sigma_{\hat{\beta}} = 2$), there is an 85.1\% probability that our test will correctly detect the effect and reject $H_0$. This is good power.
\end{solution}

\item Compute the \textbf{power} when $\beta^* = 2$ (the program raises wages by only \$2{,}000). Compare with part~(c) and explain.

\textcolor{hintcolor}{\emph{Hint: $\delta = 2/2 = 1$, so $Z \sim N(1,1)$. You may use: $\Phi(0.96) = 0.831$.}}

\begin{solution}
\textbf{Step 1: Distribution of $Z$ under $\beta^* = 2$.}
\[
  Z = \frac{\hat{\beta}}{2} \sim N\!\left(\frac{2}{2},\, 1\right) = N(1, 1)
\]
Now $\delta = 1$ (the shift from the null is smaller).

\textbf{Step 2: Compute power.}

With $W = Z - 1 \sim N(0,1)$:

\textbf{Right tail:}
\[
  P(Z > 1.96) = P(W > 0.96) = 1 - \Phi(0.96) = 1 - 0.831 = 0.169
\]

\textbf{Left tail:}
\[
  P(Z < -1.96) = P(W < -2.96) = \Phi(-2.96) \approx 0.002
\]

\textbf{Combining:}
\[
  \text{Power} = 0.169 + 0.002 = \boxed{0.171}
\]

\textbf{Comparison:}

\medskip
\begin{center}
\begin{tabular}{cccc}
\toprule
True $\beta^*$ & $\delta = \beta^*/\sigma_{\hat{\beta}}$ & Power & Detect? \\
\midrule
$6$ & $3$ & $0.851$ & Likely yes \\
$2$ & $1$ & $0.171$ & Likely no \\
\bottomrule
\end{tabular}
\end{center}

\medskip
\textbf{Why the power is so different:} Power depends on how far the true $\beta^*$ is from the null value ($\beta_0 = 0$), measured in units of $\sigma_{\hat{\beta}}$. This ratio $\delta = \beta^*/\sigma_{\hat{\beta}}$ is the ``signal-to-noise ratio.''

\begin{itemize}
  \item When $\beta^* = 6$: the signal ($\delta = 3$) is strong relative to the noise. The distribution of $Z$ is shifted far enough from zero that it almost always exceeds the critical value $1.96$. The test detects the effect 85\% of the time.

  \item When $\beta^* = 2$: the signal ($\delta = 1$) is weak. The distribution of $Z$ overlaps heavily with the rejection region. The test misses the effect about 83\% of the time (Type~II error probability $= 1 - 0.171 = 0.829$).
\end{itemize}

\textbf{Implication:} Small effects are hard to detect unless $\sigma_{\hat{\beta}}$ is small (which requires large $n$ or large $\text{SST}_x$). This connects back to Block~2: increasing $\text{SST}_x$ reduces $\text{Var}(\hat{\beta}_1) = \sigma^2/\text{SST}_x$, which reduces $\sigma_{\hat{\beta}}$, which increases $\delta$, which increases power.
\end{solution}

\end{enumerate}

%=============================================================================
\section*{Summary of Key Formulas}
%=============================================================================

\begin{center}
\renewcommand{\arraystretch}{1.6}
\begin{tabular}{ll}
\toprule
\textbf{Concept} & \textbf{Formula} \\
\midrule
Total variation in $x$ & $\text{SST}_x = \sum_{i=1}^n (x_i - \bar{x})^2$ \\
Variance of $\hat{\beta}_1$ (Thm.~2.2) & $\text{Var}(\hat{\beta}_1 \mid \mathbf{x}) = \sigma^2 / \text{SST}_x$ \;\; [uses SLR.1--5]\\
Estimated error variance & $\hat{\sigma}^2 = \sum \hat{u}_i^2 / (n-2)$ \\
Standard error of $\hat{\beta}_1$ & $\text{SE}(\hat{\beta}_1) = \hat{\sigma} / \sqrt{\text{SST}_x}$ \\
$t$-statistic & $t = (\hat{\beta}_1 - \beta_{1,0}) / \text{SE}(\hat{\beta}_1) \sim t_{n-2}$ \;\; [uses SLR.1--6]\\
Two-sided rejection rule & Reject if $|t| > t_{\alpha/2,\, n-2}$ \\
One-sided rejection rule & Reject if $t > t_{\alpha,\, n-2}$ \;\; (for $H_1\colon \beta_1 > \beta_{1,0}$) \\
$(1-\alpha)$ Confidence interval & $\hat{\beta}_1 \pm t_{\alpha/2,\, n-2} \cdot \text{SE}(\hat{\beta}_1)$ \\
CI--test equivalence & Reject two-sided $H_0$ iff $\beta_{1,0} \notin$ CI \\
\midrule
Type I Error (size) & $P(\text{reject } H_0 \mid H_0 \text{ true}) = \alpha$ \\
Type II Error & $P(\text{fail to reject } H_0 \mid H_0 \text{ false})$ \\
Power & $P(\text{reject } H_0 \mid H_0 \text{ false}) = 1 - P(\text{Type II})$ \\
\bottomrule
\end{tabular}
\end{center}

\end{document}
